\documentclass[a4paper,10pt]{article}
\usepackage[margin=22mm]{geometry}
\usepackage{fontspec,fvextra}
\setmainfont{STIXGeneral}\setmonofont{Menlo}[Scale=0.90]
\title{Computations for thin covers}\author{Human Compression project}
\date{Private author edition}
\begin{document}\maketitle
The first complete program is a required finite certificate identifying the
order96 source presentation with the marked affine model. The remaining
GAP and Python programs are bounded validations of written arguments,
including the thin-base counterexample and four sample moduli. They are
not premises of the uniform proof. The source refers to GAP and an Atlas
but supplies no full executable listing; its historical computational
cost is unknown. Optional validation is not compared asymmetrically
with an unknown source cost. Drivers, integrity tools and scoped Lean
are inventoried separately. No length factor is claimed here.
\section*{Required finite presentation certificate}
\VerbatimInput[fontsize=\footnotesize,breaklines=true,breakanywhere=true,breaksymbolleft={},breaksymbolright={}]{base_certificate.g}
\section*{Optional base and voltage checks}
\VerbatimInput[fontsize=\footnotesize,breaklines=true,breakanywhere=true,breaksymbolleft={},breaksymbolright={}]{validate_base.g}
\section*{Optional cover and counterexample checks}
\VerbatimInput[fontsize=\footnotesize,breaklines=true,breakanywhere=true,breaksymbolleft={},breaksymbolright={}]{validate_covers.py}
\end{document}
